\documentclass[../../数学分析培优讲义(答案版).tex]{subfiles}

\begin{document}

\setcounter{chapter}{12}
\pagestyle{fancy}

\chapter{曲线积分与曲面积分}

\section{曲线积分}

\subsection{第一型曲线积分}

\begin{example}
设
\[
    L:\ x^2+y^2+z^2=a^2,\qquad x+y+z=0,
\]
计算积分:
\begin{enumerate}[label=(\arabic*)]
    \item $\displaystyle\int\limits_L x^2\,\mathrm{d}s$;
    \item $\displaystyle\int\limits_L(xy+yz+zx)\,\mathrm{d}s$.
\end{enumerate}
\end{example}

\begin{proof}[解]
曲线 $L$ 是半径为 $a$ 的大圆，长度为 $2\pi a$。由坐标置换对称性，
\[
\int_Lx^2\,\mathrm{d}s=\dfrac{1}{3}\int_L(x^2+y^2+z^2)\,\mathrm{d}s=\dfrac{2\pi a^3}{3}.
\]
又因 $x+y+z=0$，有 $xy+yz+zx=-a^2/2$，故第二个积分为 $-\pi a^3$。
\end{proof}

\begin{example}
计算下列第一型曲线积分:
\begin{enumerate}[label=(\arabic*)]
    \item $\displaystyle I=\int\limits_L\sqrt{x^2+y^2}\,\mathrm{d}s$, 其中 $L:x^2+y^2=ax$ $(a>0)$;
    \item $\displaystyle I=\int\limits_L\left(x^{\frac{4}{3}}+y^{\frac{4}{3}}\right)\,\mathrm{d}s$, 其中 $L:x^{\frac{2}{3}}+y^{\frac{2}{3}}=a^{\frac{2}{3}}$ $(a>0)$.
\end{enumerate}
\end{example}

\begin{proof}[解]
\begin{enumerate}[label=(\arabic*)]
\item 参数化圆 $x=\dfrac{a}{2}(1+\cos t),y=\dfrac{a}{2}\sin t$，得 $I=2a^2$。
\item 参数化星形线 $x=a\cos^3t,y=a\sin^3t$。利用四象限对称性，
\[
I=12a^{\frac{7}{3}}\int_0^{\frac{\pi}{2}}(\cos^4t+\sin^4t)\sin t\cos t\,\mathrm{d}t=4a^{\frac{7}{3}}.
\]
\end{enumerate}
\end{proof}

\begin{example}
计算
\[
    \int\limits_L(x+y)\,\mathrm{d}s,
\]
其中 $L$ 为顶点 $(0,0),(1,0),(0,1)$ 的三角形边界.
\end{example}

\begin{proof}[解]
三条边分别积分。两条直角边各给出 $1/2$，斜边上 $x+y=1$ 且长度为 $\sqrt{2}$，故结果为 $1+\sqrt{2}$。
\end{proof}

\subsection{第二型曲线积分}

\begin{example}
求曲线积分
\[
    \int\limits_C e^{-(x^2+y^2)}\bigl\lbrack \cos(2xy)\,\mathrm{d}x+\sin(2xy)\,\mathrm{d}y\bigr\rbrack,
\]
其中 $C:x^2+y^2=1$, 方向为逆时针.
\end{example}

\begin{proof}[解]
令 $x=\cos t,y=\sin t$。积分化为
\[
e^{-1}\int_0^{2\pi}\bigl[-\sin t\cos(\sin2t)+\cos t\sin(\sin2t)\bigr]\,\mathrm{d}t.
\]
两项分别关于适当的半周期变换成相反数，故积分为 $0$。
\end{proof}

\begin{example}
利用对称性证明
\[
    I=\int\limits_L\dfrac{-x\,\mathrm{d}x+y\,\mathrm{d}y}{x^2+y^2}=0,
\]
其中 $L:x^2+y^2=1$, 方向为逆时针.
\end{example}

\begin{proof}[解]
令 $x=\cos t,y=\sin t$，则被积式为 $\sin2t\,\mathrm{d}t$，从 $0$ 到 $2\pi$ 的积分为 $0$。
\end{proof}

\begin{example}
计算积分
\[
    \int\limits_{ABC}\dfrac{\mathrm{d}x+\mathrm{d}y}{|x|+|y|},
\]
其中 $ABC$ 为三点 $A(1,0),B(0,1),C(-1,0)$ 连成的折线.
\end{example}

\begin{proof}[解]
在线段 $AB$ 上 $|x|+|y|=1$ 且 $\mathrm{d}x+\mathrm{d}y=0$；在线段 $BC$ 上仍有 $|x|+|y|=1$，而 $\mathrm{d}x+\mathrm{d}y=-2\,\mathrm{d}t$。故积分为 $-2$。
\end{proof}

\begin{example}
计算
\[
    \int\limits_L y\,\mathrm{d}x+z\,\mathrm{d}y+x\,\mathrm{d}z,
\]
其中
\[
    L:\ \dfrac{x^2}{a^2}+\dfrac{y^2}{b^2}+\dfrac{z^2}{c^2}=1,
    \qquad \dfrac{x}{a}+\dfrac{z}{c}=1,
    \qquad x,y,z\geq0,
\]
积分方向从 $(a,0,0)$ 到 $(0,0,c)$.
\end{example}

\begin{proof}[解]
令
\[
x=a\cos^2t,\quad y=\sqrt{2}b\sin t\cos t,\quad z=c\sin^2t,\qquad 0\leq t\leq\dfrac{\pi}{2}.
\]
代入并逐项积分，得到
\[
\int_Ly\,\mathrm{d}x+z\,\mathrm{d}y+x\,\mathrm{d}z
=\dfrac{ac}{2}-\dfrac{\sqrt{2}\pi b(a+c)}{8}.
\]
\end{proof}

\begin{example}
计算第一型曲线积分
\[
    \int\limits_C\left(
    \dfrac{xy}{ab}+\dfrac{\sqrt{2}yz}{b\sqrt{a^2+b^2}}+
    \dfrac{\sqrt{2}zx}{a\sqrt{a^2+b^2}}
    \right)\,\mathrm{d}s,
\]
其中 $C$ 为
\[
    \dfrac{x^2}{a^2}+\dfrac{y^2}{b^2}+\dfrac{2z^2}{a^2+b^2}=1,
    \qquad \dfrac{x}{a}+\dfrac{y}{b}=1,
    \qquad x,y,z>0
\]
的交线.
\end{example}

\begin{proof}[解]
取参数
\[
x=a\cos^2t,\quad y=b\sin^2t,\quad z=\sqrt{a^2+b^2}\sin t\cos t,\quad 0<t<\dfrac{\pi}{2}.
\]
此时 $\mathrm{d}s=\sqrt{a^2+b^2}\,\mathrm{d}t$，直接积分得
\[
I=\dfrac{\pi+8\sqrt{2}}{16}\sqrt{a^2+b^2}.
\]
\end{proof}

\begin{example}
证明
\[
    \lim\limits_{R\to+\infty}
    \int\limits_{x^2+y^2=R^2}
    \dfrac{y\,\mathrm{d}x-x\,\mathrm{d}y}{|x^2+xy+y^2|^2}=0.
\]
\end{example}

\begin{proof}[解]
二次型 $x^2+xy+y^2$ 正定，故在圆周上其绝对值不小于 $cR^2$，其中 $c>0$ 与 $R$ 无关；而 $|y\,\mathrm{d}x-x\,\mathrm{d}y|=R^2\,\mathrm{d}t$。因此积分绝对值为 $O(R^{-2})$，极限为 $0$。
\end{proof}

\begin{example}
计算
\[
    \int\limits_{\widehat{OA}}(x^2+y^2)\,\mathrm{d}x+4xy\,\mathrm{d}y,
\]
其中:
\begin{enumerate}[label=(\arabic*)]
    \item $\widehat{OA}$ 为上半圆周 $x^2+y^2=ax$ $(y\geq0)$;
    \item $\widehat{OA}$ 为 $x$ 轴上线段 $OA$.
\end{enumerate}
\end{example}

\begin{proof}[解]
记 $O=(0,0),A=(a,0)$，方向均为 $O$ 到 $A$。
\begin{enumerate}[label=(\arabic*)]
\item 将上半圆弧与线段 $AO$ 合成顺时针闭曲线，用 Green 公式得弧上积分为 $a^3/6$。
\item 在线段 $OA$ 上 $y=0$，故积分为 $\int_0^ax^2\,\mathrm{d}x=a^3/3$。
\end{enumerate}
\end{proof}

\begin{example}
计算
\[
    I=\int\limits_L(y-z)\,\mathrm{d}x+(z-x)\,\mathrm{d}y+(x-y)\,\mathrm{d}z,
\]
其中 $L:x^2+y^2+z^2=a^2$, $x+y+z=0$, 从 $z$ 轴正向看 $L$ 取逆时针方向.
\end{example}

\begin{proof}[解]
取与题给方向相容的单位法向量 $\bm n=(1,1,1)/\sqrt{3}$。向量场的旋度为 $-2(1,1,1)$，而圆盘面积为 $\pi a^2$。由 Stokes 公式，
\[
I=-2(1,1,1)\cdot\bm n\,\pi a^2=-2\sqrt{3}\pi a^2.
\]
\end{proof}

\begin{example}
计算
\[
    \int\limits_{L^+}x\,\mathrm{d}y-y\,\mathrm{d}x,
\]
其中
\[
    L^+:x^{2n+1}+y^{2n+1}=ax^ny^n,\qquad x\geq0,\ y\geq0,
\]
取逆时针方向.
\end{example}

\begin{proof}[解]
由 Green 公式，所求积分是区域面积的两倍。极坐标下边界为
\[
r=\dfrac{a\cos^n\theta\sin^n\theta}{\cos^{2n+1}\theta+\sin^{2n+1}\theta}.
\]
令 $t=\tan\theta$，则
\[
I=\int_0^{\frac{\pi}{2}}r^2\,\mathrm{d}\theta
=a^2\int_0^{+\infty}\dfrac{t^{2n}}{(1+t^{2n+1})^2}\,\mathrm{d}t
=\dfrac{a^2}{2n+1}.
\]
\end{proof}

\subsection{Green 公式}

\begin{example}
计算下列积分:
\begin{enumerate}[label=(\arabic*)]
    \item $\displaystyle\int\limits_\Gamma xy^2\,\mathrm{d}y-x^2y\,\mathrm{d}x$, $\Gamma:x^2+y^2=a^2$, 按逆时针方向;
    \item $\displaystyle\int\limits_\Gamma(x+y)\,\mathrm{d}x-(x-y)\,\mathrm{d}y$, $\Gamma:\dfrac{x^2}{a^2}+\dfrac{y^2}{b^2}=1$, 按逆时针方向.
\end{enumerate}
\end{example}

\begin{proof}[解]
用 Green 公式。
\begin{enumerate}[label=(\arabic*)]
\item 面积分为 $\iint_{x^2+y^2\leq a^2}(x^2+y^2)\,\mathrm{d}x\,\mathrm{d}y=\pi a^4/2$。
\item 旋度为 $-2$，椭圆面积为 $\pi ab$，故积分为 $-2\pi ab$。
\end{enumerate}
\end{proof}

\begin{example}
计算:
\begin{enumerate}[label=(\arabic*)]
    \item $\displaystyle\int\limits_\Gamma\dfrac{x\,\mathrm{d}y-y\,\mathrm{d}x}{4x^2+y^2}$, $\Gamma:(x-1)^2+y^2=R^2$ $(R>1)$, 按逆时针方向;
    \item $\displaystyle\int\limits_L\dfrac{x\,\mathrm{d}y-y\,\mathrm{d}x}{2(Ax^2+2Bxy+Cy^2)}$, $A,C>0$, $AC-B^2>0$, $L:x^2+y^2=1$, 按逆时针方向;
    \item $\displaystyle\int\limits_L\dfrac{x\,\mathrm{d}y-y\,\mathrm{d}x}{(ax+by)^2+(cx+dy)^2}$, $ad-bc\ne0$, $L:(ax+by)^2+(cx+dy)^2=1$, 按逆时针方向.
\end{enumerate}
\end{example}

\begin{proof}[解]
\begin{enumerate}[label=(\arabic*)]
\item 令 $u=2x,v=y$，被积式为 $\dfrac{1}{2}\mathrm{d}\arg(u+iv)$。曲线绕原点一次，故结果为 $\pi$。
\item 令 $x=\cos t,y=\sin t$，利用正定二次型的标准积分，结果为
\[
\dfrac{1}{2}\int_0^{2\pi}\dfrac{\mathrm{d}t}{A\cos^2t+2B\sin t\cos t+C\sin^2t}
=\dfrac{\pi}{\sqrt{AC-B^2}}.
\]
\item 令 $u=ax+by,v=cx+dy$。考虑线性变换的定向后，结果为
\[
\dfrac{2\pi}{|ad-bc|}.
\]
\end{enumerate}
\end{proof}

\begin{example}
计算
\[
    \int\limits_C\dfrac{e^y}{x^2+y^2}
    \bigl\lbrack (x\sin x+y\cos x)\,\mathrm{d}x+(y\sin x-x\cos x)\,\mathrm{d}y\bigr\rbrack,
\]
其中 $C:x^2+y^2=1$, 按逆时针方向.
\end{example}

\begin{proof}[解]
在单位圆上令 $x=\cos t,y=\sin t$。径向部分为零，原积分化为
\[
-\int_0^{2\pi}e^{\sin t}\cos(\cos t)\,\mathrm{d}t.
\]
它是 $-\operatorname{Re}\int_0^{2\pi}e^{i e^{-it}}\,\mathrm{d}t=-2\pi$。
\end{proof}

\begin{example}
设 $C$ 为对称于坐标轴的光滑曲线, 证明
\[
    \int\limits_C(x^3y+e^y)\,\mathrm{d}x+(xy^3+xe^y-\cos x)\,\mathrm{d}y=0.
\]
\end{example}

\begin{proof}[解]
由 Green 公式，积分等于
\[
\iint_D(y^3-x^3+\sin x)\,\mathrm{d}x\,\mathrm{d}y.
\]
区域关于两坐标轴对称，三个被积项均为相应的奇函数，故积分为 $0$。
\end{proof}

\begin{example}
设 $D$ 为第一象限内由 $y=x$, $y=4x$, $xy=1$, $xy=4$ 所围区域, $F'(x)=f(x)$, 证明
\[
    \int\limits_{\partial D}\dfrac{F(xy)}{y}\,\mathrm{d}y
    =\ln2\int_1^4f(u)\,\mathrm{d}u,
\]
其中 $\partial D$ 取逆时针方向.
\end{example}

\begin{proof}[解]
Green 公式给出左端为 $\iint_Df(xy)\,\mathrm{d}x\,\mathrm{d}y$。令 $u=xy,v=y/x$，则 $1\leq u,v\leq4$，Jacobian 为 $1/(2v)$，所以
\[
\iint_Df(xy)\,\mathrm{d}x\,\mathrm{d}y
=\dfrac{1}{2}\int_1^4f(u)\,\mathrm{d}u\int_1^4\dfrac{\mathrm{d}v}{v}
=\ln2\int_1^4f(u)\,\mathrm{d}u.
\]
\end{proof}

\begin{example}
计算开口弧段上的曲线积分:
\begin{enumerate}[label=(\arabic*)]
    \item $\displaystyle\int\limits_{C^+}(-2xe^{-x^2}\sin y)\,\mathrm{d}x+(e^{-x^2}\cos y+x^4)\,\mathrm{d}y$, 其中 $C^+$ 为从 $(1,0)$ 到 $(-1,0)$ 的上半圆 $y=\sqrt{1-x^2}$;
    \item $\displaystyle\int\limits_{\widehat{OA}}(y^2-\cos y)\,\mathrm{d}x+x\sin y\,\mathrm{d}y$, 其中 $\widehat{OA}$ 为从 $O(0,0)$ 到 $A(\pi,0)$ 的弧 $y=\sin x$.
\end{enumerate}
\end{example}

\begin{proof}[解]
\begin{enumerate}[label=(\arabic*)]
\item 前两项为 $\mathrm{d}(e^{-x^2}\sin y)$。余项在参数 $x=\cos t,y=\sin t$ 下为 $\cos^5t\,\mathrm{d}t$，故积分为 $0$。
\item 直接参数化 $y=\sin x$，并对 $x\sin y\cos x$ 分部积分，所有非初等项抵消，结果为 $-\pi/2$。
\end{enumerate}
\end{proof}

\begin{example}
设 $L$ 为平面上逐段光滑闭曲线, $D$ 为 $L$ 所围区域, $u=u(x,y)$ 在 $D$ 内直到边界有二阶连续偏导数, 试证
\[
    \int\limits_L\dfrac{\partial u}{\partial\bm n}\,\mathrm{d}s
    =\iint\limits_D\left(\dfrac{\partial^2u}{\partial x^2}+\dfrac{\partial^2u}{\partial y^2}\right)\,\mathrm{d}x\,\mathrm{d}y,
\]
其中 $\bm n$ 为 $L$ 的外法向量.
\end{example}

\begin{proof}[解]
正向边界上 $\bm n\,\mathrm{d}s=(\mathrm{d}y,-\mathrm{d}x)$。因此
\[
\int_Lu_{\bm n}\,\mathrm{d}s=\int_Lu_x\,\mathrm{d}y-u_y\,\mathrm{d}x
=\iint_D(u_{xx}+u_{yy})\,\mathrm{d}x\,\mathrm{d}y
\]
即为 Green 公式。
\end{proof}

\subsection{理论证明与应用}

\begin{example}
举例说明对第二型曲线积分, ``积分中值定理'' 不再成立.
\end{example}

\begin{proof}[解]
取单位圆正向一周并令被积函数为 $-y\,\mathrm{d}x$。此时 $\int_C\mathrm{d}x=0$，但 $\int_C-y\,\mathrm{d}x=\pi$，不可能写成某个函数中值乘以 $\int_C\mathrm{d}x$。
\end{proof}

\begin{example}
设 $P,Q,R$ 在 $L$ 上连续, $L$ 为光滑弧段, 弧长为 $l$, 证明
\[
    \left|\int\limits_LP\,\mathrm{d}x+Q\,\mathrm{d}y+R\,\mathrm{d}z\right|\leq Ml,
\]
其中
\[
    M=\max_{(x,y,z)\in L}\sqrt{P^2+Q^2+R^2}.
\]
\end{example}

\begin{proof}[解]
用弧长 $s$ 参数化，切向量长度为 $1$。由 Cauchy--Schwarz 不等式，
\[
|P x'(s)+Q y'(s)+R z'(s)|\leq\sqrt{P^2+Q^2+R^2}\leq M.
\]
在 $0\leq s\leq l$ 上积分即得结论。
\end{proof}

\begin{example}
证明
\[
    \int\limits_L\cos(\bm l,\bm n)\,\mathrm{d}s=0,
\]
其中 $L$ 为逐段光滑封闭曲线, $\bm l$ 为任意给定方向, $\bm n$ 是 $L$ 的外法线方向.
\end{example}

\begin{proof}[解]
该积分是常向量 $\bm l$ 通过闭曲线的外向通量。其散度为零，故由 Green 公式积分为 $0$。
\end{proof}

\begin{example}
计算 Gauss 积分
\[
    \int\limits_L\dfrac{\cos(\bm r,\bm n)}{r}\,\mathrm{d}s,
\]
其中 $r=\sqrt{(\xi-x)^2+(\eta-y)^2}$, $\bm r$ 为从点 $A(x,y)$ 指向封闭光滑曲线 $L$ 上动点 $M(\xi,\eta)$ 的向量.
\end{example}

\begin{proof}[解]
这是向量场 $\bm r/r^2$ 的外向通量。若 $A$ 在 $L$ 内，挖去以 $A$ 为心的小圆并用 Green 公式，结果为 $2\pi$；若 $A$ 在外部，结果为 $0$。若 $A\in L$，原积分含奇点而题面未定义。
\end{proof}

\begin{example}[\markstar]
设 $\bm F=P\bm i+Q\bm j$ 在开区域 $D$ 内处处连续可微, 在 $D$ 内任一圆周 $C$ 上有
\[
    \int\limits_C\bm F\cdot\bm n\,\mathrm{d}s=0,
\]
其中 $\bm n$ 为圆周的单位外法向量, 试证
\[
    \dfrac{\partial P}{\partial x}+\dfrac{\partial Q}{\partial y}=0
\]
在 $D$ 内恒成立.
\end{example}

\begin{proof}[解]
由 Green 公式，对任意闭圆盘 $B\Subset D$ 有
\[
0=\int_{\partial B}\bm F\cdot\bm n\,\mathrm{d}s
=\iint_B(P_x+Q_y)\,\mathrm{d}x\,\mathrm{d}y.
\]
令圆盘半径趋于零并利用连续性，便得 $P_x+Q_y=0$。
\end{proof}

\begin{example}[\markstar]
设 $P,Q$ 在全平面上有连续偏导数, 且以任意点 $(x_0,y_0)$ 为中心、任意 $r>0$ 为半径的上半圆 $C$ 上恒有
\[
    \int\limits_C P(x,y)\,\mathrm{d}x+Q(x,y)\,\mathrm{d}y=0,
\]
证明
\[
    P(x,y)\equiv0,
    \qquad
    \dfrac{\partial Q}{\partial y}\equiv0.
\]
\end{example}

\begin{proof}[解]
按通常的“上半圆弧”参数化 $x=x_0+r\cos t,y=y_0+r\sin t$，$0\leq t\leq\pi$。除以 $r$ 后令 $r\to0$，得到 $P(x_0,y_0)=0$。再展开到下一阶则得到 $Q_x(x_0,y_0)=0$。

题面所写 $Q_y\equiv0$ 不成立。例如 $P=0,Q=y$ 时任意上半圆弧上 $\int Q\,\mathrm{d}y=\frac{1}{2}[y^2]_{y_0}^{y_0}=0$，但 $Q_y=1$。因此原结论应核对，按现有题面正确条件是 $P\equiv0,Q_x\equiv0$。
\end{proof}

\begin{example}
计算
\[
    \int\limits_{x^2+y^2=R^2}\ln\sqrt{(x-a)^2+(y-b)^2}\,\mathrm{d}s.
\]
\end{example}

\begin{proof}[解]
记 $\rho=\sqrt{a^2+b^2}$。圆周上对数势的平均值给出
\[
\int_{x^2+y^2=R^2}\ln\sqrt{(x-a)^2+(y-b)^2}\,\mathrm{d}s
=2\pi R\ln\max\{R,\rho\}.
\]
当 $\rho=R$ 时按反常积分理解，仍取同一值。
\end{proof}

\begin{example}
设 $D:x^2+y^2\leq a^2$ $(a>0)$, $f(x,y)$ 在 $D$ 上有连续偏导数, 且 $f=0$ 在 $\partial D$ 上成立, 证明
\[
    \left|\iint\limits_Df(x,y)\,\mathrm{d}x\,\mathrm{d}y\right|
    \leq\dfrac{\pi a^3}{3}\max_{(x,y)\in D}
    \sqrt{\left(\dfrac{\partial f}{\partial x}\right)^2+
    \left(\dfrac{\partial f}{\partial y}\right)^2}.
\]
\end{example}

\begin{proof}[解]
记梯度最大值为 $M$。沿半径到边界应用中值定理得 $|f(r,\theta)|\leq M(a-r)$，从而
\[
\left|\iint_Df\right|\leq2\pi M\int_0^a(a-r)r\,\mathrm{d}r=\dfrac{\pi a^3}{3}M.
\]
\end{proof}

\methodsection{(A) 调和函数 *}

\begin{example}
设 $L$ 为平面上一条自身不相交的光滑曲线, 起点 $(1,0)$, 终点 $(0,2)$, 除起终点外全部位于第一象限, 计算
\[
    \int\limits_L\dfrac{\partial\ln r}{\partial\bm n}\,\mathrm{d}s,
\]
其中 $r$ 表示 $L$ 上动点到原点的距离.
\end{example}

\begin{proof}[解]
开曲线的法向必须先约定。若按起点到终点的方向取左法向，则
\[
\dfrac{\partial\ln r}{\partial\bm n}\,\mathrm{d}s
=\dfrac{y\,\mathrm{d}x-x\,\mathrm{d}y}{x^2+y^2}=-\mathrm{d}\arg(x+iy),
\]
故积分为 $-\pi/2$；取右法向则为 $\pi/2$。原题未指定法向，答案相差一个符号。
\end{proof}

\begin{example}
设 $\Omega$ 为 $xOy$ 平面上具有光滑边界的有界闭区域, $u$ 在 $\Omega$ 内有二阶连续偏导数, 直到边界有一阶连续偏导数, 且 $u|_{\partial\Omega}=0$, $u$ 非常值, 证明
\[
    \iint\limits_\Omega u\left(\dfrac{\partial^2u}{\partial x^2}+\dfrac{\partial^2u}{\partial y^2}\right)\,\mathrm{d}x\,\mathrm{d}y<0.
\]
\end{example}

\begin{proof}[解]
Green 第一恒等式给出
\[
\iint_\Omega u\Delta u\,\mathrm{d}A
=\int_{\partial\Omega}u\dfrac{\partial u}{\partial\bm n}\,\mathrm{d}s
-\iint_\Omega|\nabla u|^2\,\mathrm{d}A
=-\iint_\Omega|\nabla u|^2\,\mathrm{d}A<0,
\]
最后的严格不等号来自 $u$ 非常值。
\end{proof}

\begin{example}
设 $u(x,y)$ 有二阶连续偏导数, 证明: $u$ 为调和函数当且仅当对任意逐段光滑闭曲线 $C$,
\[
    \int\limits_C\dfrac{\partial u}{\partial\bm n}\,\mathrm{d}s=0,
\]
其中 $\bm n$ 为外法线方向.
\end{example}

\begin{proof}[解]
若 $\Delta u=0$，由上一节的通量公式，任意闭曲线上的法向导数积分为零。反之，对任意小圆盘积分均为零，除以圆盘面积并令半径趋于零，得到 $\Delta u=0$。
\end{proof}

\begin{example}
设 $D$ 为区域, $u\in C^{(2)}(D)$, 证明: $u$ 为 $D$ 上调和函数当且仅当对任意 $P_0(x_0,y_0)\in D$,
\[
    u(x_0,y_0)=\dfrac{1}{2\pi}\int_0^{2\pi}u(x_0+r\cos\theta,y_0+r\sin\theta)\,\mathrm{d}\theta,
\]
其中 $0<r<d(P_0,\partial D)$.
\end{example}

\begin{proof}[解]
固定 $P_0$，记圆周平均值为 $M(r)$。对 $r$ 求导并用极坐标及通量公式可得
\[
M'(r)=\dfrac{1}{2\pi r}\iint_{B(P_0,r)}\Delta u\,\mathrm{d}A.
\]
故调和时 $M(r)$ 为常数且极限为 $u(P_0)$。反之，若平均值性质对每个小圆成立，则上式及连续性推出 $\Delta u(P_0)=0$。
\end{proof}

\begin{example}
设 $L$ 是 $\mathbb{R}^2$ 中简单闭曲线, 其所围区域为 $D$, $f$ 在 $D$ 上连续且为调和函数, 证明 $f$ 在 $L$ 上取到最大值和最小值.
\end{example}

\begin{proof}[解]
连续函数在紧集 $\overline D$ 上取到极值。若非常值调和函数在内部取到最大值或最小值，平均值性质迫使其在一个邻域、继而在整个连通区域为常值，矛盾。因此极值均可在边界 $L$ 上取到。
\end{proof}

\begin{example}
设闭区域 $D$ 由有限条逐段光滑曲线围成, 调和函数 $u(x,y)$ 在边界 $\partial D$ 上为零, 证明 $u$ 在 $D$ 上恒为零.
\end{example}

\begin{proof}[解]
由最大值原理，$u$ 在 $D$ 上的最大值不超过边界最大值 $0$；对 $-u$ 再用一次，得 $u\geq0$。故 $u\equiv0$。
\end{proof}

\begin{example}
设 $f(x,y,z)$ 为 $n$ 次齐次函数, $f\in C^{(2)}(\mathbb{R}^3)$, 若对任意球面 $S$ 都有
\[
    \iint\limits_Sf(x,y,z)\,\mathrm{d}S=0,
\]
证明 $\Delta f=0$.
\end{example}

\begin{proof}[解]
若“任意球面”指任意球心和任意充分小半径，则球面平均恒为零。令半径趋于零，由连续性得 $f$ 在每个球心处都为零，故 $f\equiv0$，特别地 $\Delta f=0$。

若原意仅指以原点为心的球面，则结论一般不成立；齐次函数在单位球面上的平均为零并不足以推出调和性。
\end{proof}

\begin{example}
设 $\Sigma$ 是 $\mathbb{R}^3$ 第一卦限内任意光滑闭曲面, 所围立体为 $\Omega$, $f(x,y,z)$ 在第一卦限上具有一阶偏导数, 给出
\[
    \iint\limits_\Sigma f(x,y,z)(x\,\mathrm{d}y\,\mathrm{d}z+y\,\mathrm{d}z\,\mathrm{d}x+z\,\mathrm{d}x\,\mathrm{d}y)=0
\]
的充要条件.
\end{example}

\begin{proof}[解]
由 Gauss 公式，题给积分等于
\[
\iiint_\Omega\bigl(xf_x+yf_y+zf_z+3f\bigr)\,\mathrm{d}V.
\]
它对第一卦限内任意光滑闭区域都为零的充要条件是
\[
xf_x+yf_y+zf_z+3f=0.
\]
这也就是 $f$ 为 $-3$ 次齐次函数的 Euler 方程；若还要求在原点连续，则只能有 $f\equiv0$。
\end{proof}

\methodsection{(B) 曲线积分与路径无关性}

\begin{example}
计算
\[
    \int\limits_{\widehat{AB}}(x^2+2xy-y^2)\,\mathrm{d}x+(x^2-2xy-y^2)\,\mathrm{d}y,
\]
其中 $A=(0,0)$, $B=(2,1)$, $\widehat{AB}$ 为任意连接 $A,B$ 的光滑曲线.
\end{example}

\begin{proof}[解]
有 $P_y=Q_x=2x-2y$，故积分与路径无关。一个势函数为
\[
\Phi=\dfrac{x^3}{3}+x^2y-xy^2-\dfrac{y^3}{3}.
\]
所以积分为 $\Phi(2,1)-\Phi(0,0)=13/3$。
\end{proof}

\begin{example}
设 $f(x)$ 在 $(-\infty,+\infty)$ 上有连续导函数, 求
\[
    \int\limits_L
    \dfrac{1+y^2f(x,y)}{y}\,\mathrm{d}x
    +\dfrac{x}{y^2}\bigl\lbrack y^2f(x,y)-1\bigr\rbrack\,\mathrm{d}y,
\]
其中 $L$ 是从 $A(3,\dfrac{2}{3})$ 到 $B(1,2)$ 的直线段.
\end{example}

\begin{proof}[解]
题面首句的 $f(x)$ 与被积式中的 $f(x,y)$ 不一致。若按结构理解为 $f(xy)=F'(xy)$，则被积式为
\[
\mathrm{d}\left(\dfrac{x}{y}+F(xy)\right).
\]
两端点均满足 $xy=2$，故积分为 $1/2-9/2=-4$。
\end{proof}

\begin{example}
求
\[
    \int\limits_L[ e^x\sin y-3(x+y)]\,\mathrm{d}x+(e^x\cos y-ax)\,\mathrm{d}y,
    \qquad a>0,
\]
其中 $L$ 从 $A=(2a,0)$ 沿 $y=\sqrt{2ax-x^2}$ 到 $B=(0,0)$.
\end{example}

\begin{proof}[解]
前两项 $e^x\sin y\,\mathrm{d}x+e^x\cos y\,\mathrm{d}y$ 是全微分。令
\[
x=a(1+\cos t),\qquad y=a\sin t,qquad0\leq t\leq\pi,
\]
代入其余两项，得到
\[
I=\dfrac{a^2}{2}\bigl(12+3\pi-a\pi\bigr).
\]
\end{proof}

\begin{example}
计算
\[
    \int\limits_Lx\ln(x^2+y^2-1)\,\mathrm{d}x+y\ln(x^2+y^2-1)\,\mathrm{d}y,
\]
其中 $L$ 是从 $A(2,0)$ 到 $B(0,2)$ 的逐段光滑曲线.
\end{example}

\begin{proof}[解]
令 $u=x^2+y^2-1$，则被积式为 $\frac{1}{2}\ln u\,\mathrm{d}u$。两端点都有 $u=3$，故积分为 $0$。
\end{proof}

\begin{example}
计算
\[
    I=\int\limits_{L^+}\dfrac{(x+y)\,\mathrm{d}x-(x-y)\,\mathrm{d}y}{x^2+y^2},
\]
其中 $L^+$ 是从 $A(-1,0)$ 到 $B(1,0)$ 的一条不通过原点的光滑曲线, 满足 $y=f(x)$, $-1\leq x\leq1$.
\end{example}

\begin{proof}[解]
被积式可写为
\[
\mathrm{d}\ln\sqrt{x^2+y^2}-\mathrm{d}\arg(x+iy).
\]
两端半径相同。若图像从原点上方通过，则辐角改变量为 $-\pi$，$I=\pi$；若从下方通过，则辐角改变量为 $\pi$，$I=-\pi$。题面未规定曲线位于哪一侧，故有这两个可能值。
\end{proof}

\begin{example}
计算下列曲线积分:
\begin{enumerate}[label=(\arabic*)]
    \item $\displaystyle\int\limits_{\widehat{AB}}x\,\mathrm{d}x+y\,\mathrm{d}y$, $A=(0,1)$, $B=(3,4)$;
    \item $\displaystyle\int\limits_{\widehat{AB}}\dfrac{y\,\mathrm{d}x-x\,\mathrm{d}y}{x^2}$, $A=(2,1)$, $B=(1,2)$, $\widehat{AB}$ 不穿过 $y$ 轴;
    \item $\displaystyle\int\limits_{\widehat{AB}}\dfrac{x\,\mathrm{d}x+y\,\mathrm{d}y}{\sqrt{x^2+y^2}}$, $A=(1,0)$, $B=(6,8)$, $\widehat{AB}$ 不过原点.
\end{enumerate}
\end{example}

\begin{proof}[解]
三式都是全微分：
\begin{enumerate}[label=(\arabic*)]
\item $\mathrm{d}[(x^2+y^2)/2]$，积分为 $12$；
\item $-\mathrm{d}(y/x)$，积分为 $-3/2$；
\item $\mathrm{d}\sqrt{x^2+y^2}$，积分为 $9$。
\end{enumerate}
\end{proof}

\begin{example}
求下列表达式的原函数:
\begin{enumerate}[label=(\arabic*)]
    \item $\displaystyle\dfrac{y\,\mathrm{d}x-x\,\mathrm{d}y}{3x^2-2xy+3y^2}$, $y>0$;
    \item $\displaystyle\dfrac{(x^2+2xy+5y^2)\,\mathrm{d}x+(x^2-2xy+y^2)\,\mathrm{d}y}{(x+y)^2}$.
\end{enumerate}
\end{example}

\begin{proof}[解]
\begin{enumerate}[label=(\arabic*)]
\item 令 $t=x/y$，原式化为 $\mathrm{d}t/(3t^2-2t+3)$，故一个原函数为
\[
\dfrac{1}{2\sqrt{2}}\arctan\dfrac{3x-y}{2\sqrt{2}y}.
\]
\item 该微分式并非恰当微分，因为直接计算得 $P_y-Q_x=4y/(x+y)^2$。所以在一般区域内不存在原函数，题面应有抄录错误。
\end{enumerate}
\end{proof}

\begin{example}
求 $P,Q\in C^{(1)}(\mathbb{R}^2)$ 的关系, 使得对任意光滑闭曲线 $L\subset\mathbb{R}^2$,
\[
    I(\alpha,\beta)=\int\limits_LP(x+\alpha,y+\beta)\,\mathrm{d}x+Q(x+\alpha,y+\beta)\,\mathrm{d}y
\]
与常数 $\alpha,\beta$ 无关.
\end{example}

\begin{proof}[解]
由 Green 公式，平移后的积分是 $Q_x-P_y$ 在平移区域上的面积分。它对任意闭曲线及任意平移均不变的充要条件是
\[
Q_x-P_y=C
\]
为常数；充分性显然，必要性可对任意小圆应用并令半径趋于零。
\end{proof}

\newpage

\section{曲面积分}

\subsection{第一型曲面积分}

\begin{example}
设 $f(z)$ 为奇函数, 求
\[
    I=\iint\limits_Sf(z)\,\mathrm{d}S,\qquad
    J=\iint\limits_Sf^2(z)\,\mathrm{d}S,\qquad
    R=\iint\limits_Syf^2(z)\,\mathrm{d}S,
\]
其中 $S$ 为锥面 $z^2=2xy$ 位于球面 $x^2+y^2+z^2=a^2$ 内的部分.
\end{example}

\begin{proof}[解]
曲面关于 $z\mapsto-z$ 及 $(x,y)\mapsto(-x,-y)$ 对称，故
\[
I=0,\qquad R=0.
\]
令 $u=(x+y)/\sqrt{2},v=(x-y)/\sqrt{2}$，以
$u=\pm r,v=r\cos\theta,z=r\sin\theta$ 参数化两叶锥面，则
\[
J=2\sqrt{2}\int_0^{a/\sqrt{2}}\!r\,\mathrm{d}r
\int_0^{2\pi}f^2(r\sin\theta)\,\mathrm{d}\theta.
\]
\end{proof}

\begin{example}
计算锥面
\[
    z=\sqrt{2xy},\qquad x\geq0,\ y\geq0,
\]
位于球面 $x^2+y^2+z^2=a^2$ 内部分的面积.
\end{example}

\begin{proof}[解]
沿用 $u=(x+y)/\sqrt{2},v=(x-y)/\sqrt{2}$。该曲面可参数化为
\[
u=r,\quad v=r\cos\theta,\quad z=r\sin\theta,\quad
0\leq r\leq\dfrac{a}{\sqrt{2}},\quad0\leq\theta\leq\pi.
\]
面积元为 $\sqrt{2}r\,\mathrm{d}r\,\mathrm{d}\theta$，故面积为 $\pi a^2/(2\sqrt{2})$。
\end{proof}

\begin{example}
计算
\[
    \iint\limits_S|z|\,\mathrm{d}S,
\]
其中 $S$ 为柱体 $x^2+y^2\leq ax$ 被球体 $x^2+y^2+z^2\leq a^2$ 截取部分的表面.
\end{example}

\begin{proof}[解]
若 $S$ 指被球截取的圆柱侧面，参数化底圆并先对 $z$ 积分可得
\[
\iint_S|z|\,\mathrm{d}S=\dfrac{\pi a^3}{2}.
\]
若“表面”指两个实体交集的整个边界，还需加上球面部分；球面部分同样为 $\pi a^3/2$，此时总值为 $\pi a^3$。原题措辞需据原始资料确定。
\end{proof}

\begin{example}
计算
\[
    \iint\limits_Sz\,\mathrm{d}S,
\]
其中 $S$ 是曲面 $x^2+z^2=2az$ $(a>0)$ 被曲面 $z=\sqrt{x^2+y^2}$ 所截取的有限部分.
\end{example}

\begin{proof}[解]
令 $x=a\sin t,z=a(1+\cos t)$。锥面条件给出
$|y|\leq a\sqrt{2\cos t(1+\cos t)}$，$|t|\leq\pi/2$，而柱面面积元为 $a\,\mathrm{d}t\,\mathrm{d}y$。因此
\[
I=4\sqrt{2}a^3\int_0^{\frac{\pi}{2}}(1+\cos t)^{\frac{3}{2}}(\cos t)^{\frac{1}{2}}\,\mathrm{d}t.
\]
按现有题面该积分一般不化为初等常数；若原资料预期简单数值，应核对柱面方程是否漏写 $y^2$。
\end{proof}

\begin{example}
计算曲面积分
\[
    F(t)=\iint\limits_{x+y+z=t}f(x,y,z)\,\mathrm{d}S,
\]
其中
\[
    f(x,y,z)=
    \begin{cases}
        1-(x^2+y^2+z^2), & x^2+y^2+z^2\leq1,\\
        0, & x^2+y^2+z^2>1.
    \end{cases}
\]
\end{example}

\begin{proof}[解]
平面到原点距离为 $|t|/\sqrt{3}$。当 $|t|\leq\sqrt{3}$ 时，截得圆盘半径平方为 $1-t^2/3$，故
\[
F(t)=\dfrac{\pi}{2}\left(1-\dfrac{t^2}{3}\right)^2.
\]
当 $|t|>\sqrt{3}$ 时平面不与单位球相交，$F(t)=0$。
\end{proof}

\begin{example}
计算
\[
    \iint\limits_\Sigma\dfrac{\mathrm{d}S}{\sqrt{x^2+y^2+(z+a)^2}},
\]
其中 $\Sigma$ 为以原点为中心、$a$ 为半径的上半球面.
\end{example}

\begin{proof}[解]
用球坐标并注意分母是上半球动点到 $(0,0,-a)$ 的距离：
\[
I=2\pi a^2\int_0^{\frac{\pi}{2}}\dfrac{\sin\varphi\,\mathrm{d}\varphi}{a\sqrt{2(1+\cos\varphi)}}
=4\pi a\left(1-\dfrac{1}{\sqrt{2}}\right).
\]
\end{proof}

\begin{example}
已知 $S=\{(x,y,z)\mid x^2+y^2+z^2=1\}$, 试用 $\Gamma$ 函数表示
\[
    \iint\limits_S(x^2+y^2)^\beta\,\mathrm{d}S,
    \qquad \beta>-1.
\]
\end{example}

\begin{proof}[解]
球坐标下 $x^2+y^2=\sin^2\varphi$，故
\[
I=2\pi\int_0^\pi\sin^{2\beta+1}\varphi\,\mathrm{d}\varphi
=\dfrac{2\pi^{3/2}\Gamma(\beta+1)}{\Gamma(\beta+\frac{3}{2})}.
\]
\end{proof}

\begin{example}
设 $f$ 连续, 证明 Poisson 公式
\[
    \int_0^{2\pi}\int_0^\pi
    f(a\sin\varphi\cos\theta+b\sin\varphi\sin\theta+c\cos\varphi)
    \sin\varphi\,\mathrm{d}\varphi\,\mathrm{d}\theta
    =2\pi\int_{-1}^1f(kz)\,\mathrm{d}z,
\]
其中 $k=\sqrt{a^2+b^2+c^2}$.
\end{example}

\begin{proof}[解]
作正交旋转，使向量 $(a,b,c)$ 变为 $(0,0,k)$。球面面积元及积分域不变，左端成为
\[
\int_0^{2\pi}\!\mathrm{d}\theta\int_0^\pi f(k\cos\varphi)\sin\varphi\,\mathrm{d}\varphi
=2\pi\int_{-1}^1f(kz)\,\mathrm{d}z.
\]
\end{proof}

\begin{example}
设 $f\in C^{(1)}\lbrack-1,1\rbrack$, $f(-1)=f(1)=0$,
\[
    M=\max_{-1\leq x\leq1}|f'(x)|,
\]
$S$ 为单位球面, 常数 $m,n,p$ 满足 $m^2+n^2+p^2=1$, 证明
\[
    \left|\iint\limits_Sf(mx+ny+pz)\,\mathrm{d}S\right|\leq2\pi M.
\]
\end{example}

\begin{proof}[解]
由上一题及旋转不变性，左端绝对值不超过 $2\pi\int_{-1}^1|f(t)|\,\mathrm{d}t$。端点为零及中值定理给出 $|f(t)|\leq M(1-|t|)$，而其积分为 $M$，故所求不超过 $2\pi M$。
\end{proof}

\begin{example}
设 $\Sigma:x^2+y^2+z^2=1$, $A$ 为 $\Sigma$ 内一点, $A$ 到原点的距离为 $q$ $(0<q<1)$, $\rho$ 表示 $A$ 到动点 $P\in\Sigma$ 的距离, 求
\[
    I=\iint\limits_\Sigma\dfrac{\mathrm{d}S}{\rho^2}.
\]
\end{example}

\begin{proof}[解]
旋转坐标使 $A=(0,0,q)$。于是
\[
I=2\pi\int_{-1}^1\dfrac{\mathrm{d}u}{1+q^2-2qu}
=\dfrac{2\pi}{q}\ln\dfrac{1+q}{1-q}.
\]
\end{proof}

\begin{example}
设 $S$ 为椭球面, $\rho$ 表示从椭球中心到椭球表面元素 $\mathrm{d}S$ 所在切平面的距离, 计算
\[
    \text{(1) }\iint\limits_S\rho\,\mathrm{d}S;
    \qquad
    \text{(2) }\iint\limits_S\dfrac{1}{\rho}\,\mathrm{d}S.
\]
\end{example}

\begin{proof}[解]
设椭球半轴为 $a,b,c$，写成 $\bm x=A\bm u$，其中 $A=\operatorname{diag}(a,b,c)$、$|\bm u|=1$。有
\[
\mathrm{d}S=abc|A^{-T}\bm u|\,\mathrm{d}\Omega,\qquad
\rho=\dfrac{1}{|A^{-T}\bm u|}.
\]
因此
\[
\iint_S\rho\,\mathrm{d}S=4\pi abc,
\]
\[
\iint_S\dfrac{1}{\rho}\,\mathrm{d}S
=\dfrac{4\pi abc}{3}\left(\dfrac{1}{a^2}+\dfrac{1}{b^2}+\dfrac{1}{c^2}\right).
\]
\end{proof}

\subsection{第二型曲面积分}

\begin{example}
设 $f$ 为奇函数, $S^+$ 为 $|x|+|y|+|z|=1$ 的外侧, 利用对称性求出或化简
\[
\begin{aligned}
    I&=\iint\limits_{S^+}(\mathrm{d}x\,\mathrm{d}y+\mathrm{d}y\,\mathrm{d}z+\mathrm{d}z\,\mathrm{d}x),\\
    J&=\iint\limits_{S^+}f^2(z)\,\mathrm{d}x\,\mathrm{d}y,\\
    K&=\iint\limits_{S^+}xf^2(z)\,\mathrm{d}x\,\mathrm{d}y,\\
    L&=\iint\limits_{S^+}f(z)\,\mathrm{d}x\,\mathrm{d}y,\\
    M&=\iint\limits_{S^+}(x+2y+3z)f(x+y+z)\,\mathrm{d}x\,\mathrm{d}y.
\end{aligned}
\]
\end{example}

\begin{proof}[解]
由对称性或 Gauss 公式，
\[
I=J=K=M=0.
\]
对 $L$ 用 Gauss 公式并按 $z$ 截面，
\[
L=\iiint_{|x|+|y|+|z|\leq1}f'(z)\,\mathrm{d}V
=2\int_{-1}^1(1-|z|)^2f'(z)\,\mathrm{d}z
=8\int_0^1(1-z)f(z)\,\mathrm{d}z.
\]
其中 $M=0$ 还用到区域对坐标置换和中心反演均不变。
\end{proof}

\begin{example}
设
\[
    S^+:z-c=\sqrt{R^2-(x-a)^2-(y-b)^2}
\]
取上侧, 计算
\[
    \iint\limits_{S^+}x^2\,\mathrm{d}y\,\mathrm{d}z+y^2\,\mathrm{d}z\,\mathrm{d}x+(x-a)y z\,\mathrm{d}x\,\mathrm{d}y.
\]
\end{example}

\begin{proof}[解]
把上半球面与底面圆盘合成闭曲面。底面通量因关于 $x-a$ 的奇对称性为零。向量场
\[
\bm F=(x^2,y^2,(x-a)yz)
\]
的散度为 $2x+2y+(x-a)y$；在以 $(a,b,c)$ 为心的上半球中，所有奇项积分为零。因此
\[
I=2(a+b)\cdot\dfrac{2\pi R^3}{3}=\dfrac{4\pi R^3(a+b)}{3}.
\]
\end{proof}

\begin{example}
计算
\[
    \dfrac{1}{3}\iint\limits_S(x\,\mathrm{d}y\,\mathrm{d}z+y\,\mathrm{d}z\,\mathrm{d}x+z\,\mathrm{d}x\,\mathrm{d}y),
\]
其中 $S$ 是球面 $x^2+y^2+z^2=a^2$ 的外侧.
\end{example}

\begin{proof}[解]
由 Gauss 公式，积分的三倍等于向量场 $(x,y,z)$ 的外通量，即 $3$ 倍球体体积。因此原积分为
\[
\dfrac{4\pi a^3}{3}.
\]
\end{proof}

\begin{example}
设 $\Delta$ 是以 $(1,0,0),(0,1,0),(0,0,1)$ 为顶点的三角形的上侧, 计算
\[
    I=\iint\limits_\Delta(x\,\mathrm{d}y\,\mathrm{d}z+y\,\mathrm{d}z\,\mathrm{d}x+z\,\mathrm{d}x\,\mathrm{d}y),
\]
\[
    J=\iint\limits_\Delta(x^2\,\mathrm{d}y\,\mathrm{d}z+y^2\,\mathrm{d}z\,\mathrm{d}x+z^2\,\mathrm{d}x\,\mathrm{d}y).
\]
\end{example}

\begin{proof}[解]
平面写成 $z=1-x-y$，上侧有
$\bm n\,\mathrm{d}S=(1,1,1)\,\mathrm{d}x\,\mathrm{d}y$，投影域为 $x,y\geq0,x+y\leq1$。故
\[
I=\iint_D1\,\mathrm{d}x\,\mathrm{d}y=\dfrac{1}{2},
\qquad
J=\iint_D\bigl[x^2+y^2+(1-x-y)^2\bigr]\,\mathrm{d}x\,\mathrm{d}y=\dfrac{1}{4}.
\]
\end{proof}

\begin{example}
求
\[
    \iint\limits_Sz^2\cos\gamma\,\mathrm{d}S,
\]
其中 $S$ 为上半球面 $x^2+y^2+z^2=1$, $z\geq0$, $\gamma$ 为球面外法线与 $z$ 轴的夹角.
\end{example}

\begin{proof}[解]
单位球面上 $\cos\gamma=z$。故
\[
I=\int_0^{2\pi}\!\mathrm{d}\theta\int_0^{\frac{\pi}{2}}\cos^3\varphi\sin\varphi\,\mathrm{d}\varphi=\dfrac{\pi}{2}.
\]
\end{proof}

\subsection{Gauss 公式和 Stokes 公式}

\begin{example}
计算下列积分:
\begin{enumerate}[label=(\arabic*)]
    \item $\displaystyle\iint\limits_\Sigma(x-y)\,\mathrm{d}y\,\mathrm{d}z+(y-z)\,\mathrm{d}z\,\mathrm{d}x+(z-x)\,\mathrm{d}x\,\mathrm{d}y$, 其中 $\Sigma:z=x^2+y^2$, $z\leq1$, 方向朝下;
    \item $\displaystyle\iint\limits_{S^+}(x+y-z)\,\mathrm{d}y\,\mathrm{d}z+[2y+\sin(z+x)]\,\mathrm{d}z\,\mathrm{d}x+(3z+e^{x+y})\,\mathrm{d}x\,\mathrm{d}y$, 其中 $S^+$ 为 $|x-y+z|+|y-z+x|+|z-x+y|=1$ 的外表面.
\end{enumerate}
\end{example}

\begin{proof}[解]
\begin{enumerate}[label=(\arabic*)]
\item 与顶圆盘合成闭曲面。散度为 $3$，体积为 $\pi/2$，顶圆盘通量为 $\pi$，故朝下抛物面上的积分为 $\pi/2$。
\item 散度为 $6$。令
$u=x-y+z,v=y-z+x,w=z-x+y$，其行列式绝对值为 $4$；标准八面体体积为 $4/3$，故原立体体积为 $1/3$，通量为 $2$。
\end{enumerate}
\end{proof}

\begin{example}
计算下列积分:
\begin{enumerate}[label=(\arabic*)]
    \item $\displaystyle\iint\limits_Sx^3\,\mathrm{d}y\,\mathrm{d}z+y^3\,\mathrm{d}z\,\mathrm{d}x+z^3\,\mathrm{d}x\,\mathrm{d}y$, 其中 $S$ 是上半椭球面 $\dfrac{x^2}{a^2}+\dfrac{y^2}{b^2}+\dfrac{z^2}{c^2}=1$, $z\geq0$, 方向朝上;
    \item $\displaystyle\iint\limits_\Sigma z\,\mathrm{d}x\,\mathrm{d}y+y\,\mathrm{d}z\,\mathrm{d}x+x\,\mathrm{d}y\,\mathrm{d}z$, 其中 $\Sigma$ 为圆柱面 $x^2+y^2=1$ 被 $z=0,z=3$ 截取部分的外侧.
\end{enumerate}
\end{example}

\begin{proof}[解]
\begin{enumerate}[label=(\arabic*)]
\item 补上底面后底面通量为零。散度为 $3(x^2+y^2+z^2)$，在上半椭球上积分得
\[
\dfrac{2\pi abc}{5}(a^2+b^2+c^2).
\]
\item 向量场为 $(x,y,z)$。圆柱侧面上外法向量为 $(x,y,0)$，故通量密度为 $1$，结果为侧面积 $6\pi$。
\end{enumerate}
\end{proof}

\begin{example}
若 $S$ 为封闭光滑曲面, $\bm l$ 为任意固定方向, 证明
\[
    \iint\limits_S\cos(\bm n,\bm l)\,\mathrm{d}S=0.
\]
\end{example}

\begin{proof}[解]
被积式是常向量 $\bm l$ 通过闭曲面的外通量。因 $\operatorname{div}\bm l=0$，Gauss 公式给出积分为 $0$。
\end{proof}

\begin{example}
记 $r=r(\theta,\varphi)$ 为分片光滑封闭曲面 $S$ 的球坐标方程, 证明 $S$ 所围区域体积满足
\[
    V=\dfrac{1}{3}\iint\limits_Sr\cos\psi\,\mathrm{d}S,
\]
其中 $\psi$ 为曲面外法线方向与径向量的夹角.
\end{example}

\begin{proof}[解]
对向量场 $\bm r=(x,y,z)$ 使用 Gauss 公式：
\[
3V=\iint_S\bm r\cdot\bm n\,\mathrm{d}S
=\iint_Sr\cos\psi\,\mathrm{d}S.
\]
两边除以 $3$ 即得。
\end{proof}

\begin{example}
设 $V$ 是不含原点的有界闭区域, 边界为光滑简单闭曲面 $\Sigma$, 单位外法向量为 $\bm n$, 径向量 $\bm r=(x,y,z)$, $f\in C^1\lbrack0,+\infty)$, 且
\[
    f'(t)+2f(t)-1=0,
\]
计算
\[
    \iint\limits_\Sigma f\left(\sqrt{x^2+y^2+z^2}\right)\cos(\bm r,\bm n)\,\mathrm{d}S.
\]
\end{example}

\begin{proof}[解]
按题面，$f'(t)+2f(t)-1=0$ 给出 $f(t)=1/2+Ce^{-2t}$。而所给曲面积分是径向场 $f(r)\bm e_r$ 的通量，Gauss 公式实际给出
\[
\iiint_V\left(f'(r)+\dfrac{2f(r)}{r}\right)\mathrm{d}V,
\]
它不能由题给常微分方程化为与区域无关的常数。若原方程应为 $f'(t)+2f(t)/t=1$，则答案才是 $V$ 的体积。故此题现有题面缺少因子 $1/t$，需要核对原始资料。
\end{proof}

\begin{example}
计算
\[
    \iint\limits_S\dfrac{\cos(\bm n,\bm r)}{r^2}\,\mathrm{d}S,
\]
其中 $S$ 是光滑封闭曲面, 原点不在 $S$ 上, $r$ 为动点到原点的距离.
\end{example}

\begin{proof}[解]
这是场 $\bm r/r^3$ 的外通量。若原点在 $S$ 内，挖去小球后用 Gauss 公式得 $4\pi$；若原点在 $S$ 外则为 $0$。
\end{proof}

\begin{example}
计算
\[
    \iint\limits_{S_{\text{外}}}\dfrac{x\,\mathrm{d}y\,\mathrm{d}z+y\,\mathrm{d}z\,\mathrm{d}x+z\,\mathrm{d}x\,\mathrm{d}y}{(x^2+y^2+z^2)^{\frac{3}{2}}},
\]
其中 $S$ 是立方体 $|x|\leq2$, $|y|\leq2$, $|z|\leq2$ 的表面.
\end{example}

\begin{proof}[解]
立方体包含原点，且取外侧，故由上一题或直接用 Gauss 公式，积分为 $4\pi$。
\end{proof}

\begin{example}
设 $V\subset\mathbb{R}^3$ 是其内任一封闭曲面均可在 $V$ 内连续收缩至一点的区域, $P,Q,R$ 在 $V$ 上有连续偏导数. 证明: 对 $V$ 内任一不自交光滑封闭曲面 $S$,
\[
    \iint\limits_S\bigl\lbrack P\cos(\bm n,x)+Q\cos(\bm n,y)+R\cos(\bm n,z)\bigr\rbrack\,\mathrm{d}S=0
\]
当且仅当
\[
    \dfrac{\partial P}{\partial x}+\dfrac{\partial Q}{\partial y}+\dfrac{\partial R}{\partial z}=0
\]
在 $V$ 内处处成立.
\end{example}

\begin{proof}[解]
若散度为零，Gauss 公式立即给出任意闭曲面的通量为零。反之，在任一点附近取充分小的球面，球内仍包含于 $V$；通量为零蕴含散度在球内的积分为零。除以体积并令半径趋于零，利用连续性即得散度处处为零。
\end{proof}

\begin{example}
设 $P(x,y,z),Q(x,y,z),R(x,y,z)$ 在 $\mathbb{R}^3$ 中有一阶偏导数. 对任意 $r>0$ 以及任意点
\[
    (x_0,y_0,z_0)\in\mathbb{R}^3,
\]
考虑半球面
\[
    S:z=z_0+\sqrt{r^2-(x-x_0)^2-(y-y_0)^2}
\]
上的积分
\[
    \iint\limits_S
    P\,\mathrm{d}y\,\mathrm{d}z
    +Q\,\mathrm{d}z\,\mathrm{d}x
    +R\,\mathrm{d}x\,\mathrm{d}y=0.
\]
证明
\[
    \dfrac{\partial P}{\partial x}+\dfrac{\partial Q}{\partial y}=0,
    \qquad R=0
\]
在 $\mathbb{R}^3$ 内处处成立.
\end{example}

\begin{proof}[解]
以任一点为球心。令半径趋于零，半球面上 $P,Q$ 的常数项通量为零，而 $R$ 的常数项通量为 $\pi r^2R(x_0,y_0,z_0)$，故 $R=0$。此后补上底面，其通量也为零；Gauss 公式给出每个半球体内 $P_x+Q_y$ 的积分为零，再令半径趋于零，得到 $P_x+Q_y=0$。
\end{proof}

\begin{example}
计算
\[
    \int\limits_{L^+}(z-y)\,\mathrm{d}x+(x-z)\,\mathrm{d}y+(y-x)\,\mathrm{d}z,
\]
其中 $L^+$ 从 $A(a,0,0)$ 经 $B(0,a,0)$ 到 $C(0,0,a)$ 再回到 $A$.
\end{example}

\begin{proof}[解]
令 $\bm F=(z-y,x-z,y-x)$，则 $\nabla\times\bm F=2(1,1,1)$。三角形 $ABC$ 的定向面积向量为
\[
\dfrac{1}{2}(B-A)\times(C-A)=\dfrac{a^2}{2}(1,1,1).
\]
由 Stokes 公式，积分为 $3a^2$。
\end{proof}

\begin{example}
计算
\[
    \int\limits_{L^+}(y^2+z^2)\,\mathrm{d}x+(x^2+z^2)\,\mathrm{d}y+(x^2+y^2)\,\mathrm{d}z,
\]
其中 $L$ 是曲面 $x^2+y^2+z^2=4x$ 与 $x^2+y^2=2x$ 的交线中 $z\geq0$ 的部分, 积分方向从原点进入第一卦限.
\end{example}

\begin{proof}[解]
参数化
\[
x=1+\cos t,\quad y=\sin t,\quad z=2\cos\dfrac{t}{2},\qquad -\pi\leq t\leq\pi.
\]
题给方向从 $t=\pi$ 向 $t=-\pi$。代入积分并化简，得到 $-4\pi$。
\end{proof}

\begin{example}
计算
\[
    \int\limits_{L^+}y\,\mathrm{d}x+z\,\mathrm{d}y+x\,\mathrm{d}z,
\]
其中 $L^+$ 为圆周 $x^2+y^2+z^2=a^2$, $x+y+z=0$, 从 $z$ 轴正方向看为逆时针.
\end{example}

\begin{proof}[解]
向量场 $(y,z,x)$ 的旋度为 $-(1,1,1)$。题给方向对应圆盘法向量 $(1,1,1)/\sqrt{3}$，圆盘面积为 $\pi a^2$，故由 Stokes 公式
\[
I=-\sqrt{3}\pi a^2.
\]
\end{proof}

\begin{example}
计算
\[
    \int\limits_{L^+}(x^2-yz)\,\mathrm{d}x+(y^2-xz)\,\mathrm{d}y+(z^2-xy)\,\mathrm{d}z,
\]
其中 $L^+$ 是从 $A(1,0,0)$ 到 $B(1,0,2)$ 的光滑曲线.
\end{example}

\begin{proof}[解]
被积式是势函数
\[
\Phi(x,y,z)=\dfrac{x^3+y^3+z^3}{3}-xyz
\]
的全微分。因此积分为
\[
\Phi(1,0,2)-\Phi(1,0,0)=\dfrac{8}{3}.
\]
\end{proof}

\begin{example}
设 $L$ 为空间某封闭光滑曲线, $P,Q,R$ 为空间中具有一阶连续偏导数的函数, 证明
\[
\left|\int\limits_{L^+}P\,\mathrm{d}x+Q\,\mathrm{d}y+R\,\mathrm{d}z\right|
\leq
\max_{(x,y,z)\in S}
\sqrt{
\left(\dfrac{\partial Q}{\partial x}-\dfrac{\partial P}{\partial y}\right)^2+
\left(\dfrac{\partial R}{\partial y}-\dfrac{\partial Q}{\partial z}\right)^2+
\left(\dfrac{\partial P}{\partial z}-\dfrac{\partial R}{\partial x}\right)^2}
\cdot S,
\]
其中 $S$ 表示以 $L$ 为边界的某曲面, 同时也用 $S$ 表示其面积.
\end{example}

\begin{proof}[解]
由 Stokes 公式，左端等于 $\iint_S(\nabla\times\bm F)\cdot\bm n\,\mathrm{d}S$。逐点使用 Cauchy--Schwarz 不等式，再以旋度模的最大值控制，便有
\[
\left|\int_L\bm F\cdot\mathrm{d}\bm r\right|
\leq\max_S|\nabla\times\bm F|\iint_S\mathrm{d}S,
\]
这正是题给不等式。
\end{proof}

\newpage

\section{场论初步 *}

\begin{example}
证明 $\nabla$ 的下列性质:
\begin{enumerate}[label=(\arabic*)]
    \item $\nabla(cf)=c\nabla f$, $c$ 为常数;
    \item $\nabla(f\pm g)=\nabla f\pm\nabla g$;
    \item $\nabla(fg)=f\nabla g+g\nabla f$;
    \item 若 $\varphi$ 为单变量函数, 则 $\nabla(\varphi\circ f)=(\varphi'\circ f)\nabla f$.
\end{enumerate}
\end{example}

\begin{proof}[解]
逐坐标求偏导即可。偏导算子的线性性给出 (1)、(2)；对每个坐标使用乘积法则得到 (3)；对每个坐标使用一元复合函数链式法则得到 (4)。
\end{proof}

\begin{example}
证明 $\operatorname{div}$ 的下列性质:
\begin{enumerate}[label=(\arabic*)]
    \item $\operatorname{div}(c\bm A)=c\operatorname{div}\bm A$, $c$ 为常数;
    \item $\operatorname{div}(\bm A\pm\bm B)=\operatorname{div}\bm A\pm\operatorname{div}\bm B$;
    \item $\operatorname{div}(u\bm A)=u\operatorname{div}\bm A+\bm A\cdot\nabla u$.
\end{enumerate}
\end{example}

\begin{proof}[解]
把 $\bm A=(P,Q,R)$ 代入 $\operatorname{div}\bm A=P_x+Q_y+R_z$。线性性立即给出 (1)、(2)，而
\[
\operatorname{div}(u\bm A)=(uP)_x+(uQ)_y+(uR)_z
=u\operatorname{div}\bm A+\bm A\cdot\nabla u
\]
给出 (3)。
\end{proof}

\begin{example}
证明 $\operatorname{rot}$ 的下列性质:
\begin{enumerate}[label=(\arabic*)]
    \item $\operatorname{rot}(c\bm A)=c\operatorname{rot}\bm A$, $c$ 为常数;
    \item $\operatorname{rot}(\bm A+\bm B)=\operatorname{rot}\bm A+\operatorname{rot}\bm B$;
    \item $\operatorname{rot}(u\bm A)=u\operatorname{rot}\bm A+\nabla u\times\bm A$;
    \item $\nabla\cdot(\bm F_1\times\bm F_2)=(\nabla\times\bm F_1)\cdot\bm F_2-(\nabla\times\bm F_2)\cdot\bm F_1$.
\end{enumerate}
\end{example}

\begin{proof}[解]
将旋度写成行列式
\[
\nabla\times\bm A=
\begin{vmatrix}\bm i&\bm j&\bm k\\
\partial_x&\partial_y&\partial_z\\P&Q&R\end{vmatrix}.
\]
线性性给出 (1)、(2)，乘积法则展开给出 (3)。对 (4) 逐分量展开两边，混合项逐一抵消，即得
$\nabla\cdot(\bm F_1\times\bm F_2)=(\nabla\times\bm F_1)\cdot\bm F_2-(\nabla\times\bm F_2)\cdot\bm F_1$。
\end{proof}

\begin{example}
若 $\operatorname{rot}\bm F=\nabla\times\bm F=0$ 在 $D$ 上处处成立, 则称 $\bm F$ 为 $D$ 上的无旋场.
\end{example}

\begin{proof}[解]
本题内容是“无旋场”的定义而不是待证命题：满足 $\nabla\times\bm F=0$ 的向量场称为无旋场，无需另作证明。
\end{proof}

\begin{example}
证明上述定理.
\end{example}

\begin{proof}[解]
这里“上述定理”指：在曲面单连通区域 $D$ 中，$C^2$ 向量场的“有势、无旋、保守”三者等价。

若 $\bm F=\nabla\varphi$，混合偏导可交换，故 $\nabla\times\bm F=0$；并且任意曲线积分只等于势函数的端点差，所以闭路积分为零。若 $\bm F$ 保守，固定 $x_0\in D$ 并定义
\[
\varphi(x)=\int_{x_0}^{x}\bm F\cdot\mathrm{d}\bm r,
\]
闭路积分为零保证定义与路径无关，沿坐标方向求导即得 $\nabla\varphi=\bm F$。最后，若 $\nabla\times\bm F=0$，曲面单连通性保证任意闭曲线在 $D$ 内张成曲面；Stokes 公式给出闭路积分为零。因此三者等价。
\end{proof}

\end{document}
